% ---------------------------------------------------------------------------
% Regularized estimation of the regression coefficients in spatio-temporal
% state-space models with spatial regimes.
%
% This is NOT part of the SC-STEM package paper. It is a self-contained draft of
% the theory behind STEM_Fit(alpha, lambda), kept in the repository because the
% material is a paper of its own and does not belong in the software paper.
%
% Compile with pdflatex; it depends only on amsmath, amssymb, booktabs, natbib.
% ---------------------------------------------------------------------------
\documentclass[11pt,a4paper]{article}
\usepackage[margin=3cm]{geometry}
\usepackage{amsmath,amssymb,amsthm}
\usepackage{booktabs}
\usepackage[round]{natbib}
\usepackage{hyperref}

\newtheorem{proposition}{Proposition}
\newtheorem{remark}{Remark}
\DeclareMathOperator{\tr}{tr}
\DeclareMathOperator*{\argmax}{arg\,max}
\DeclareMathOperator{\sgn}{sign}
\newcommand{\Sig}{\Sigma}
\newcommand{\Se}{\Sig_{e}}

\title{Regularized regression coefficients in spatio-temporal\\
       state-space models with spatial regimes}
\author{Draft --- not for circulation}
\date{\today}

\begin{document}
\maketitle

\begin{abstract}
We add an elastic-net penalty on the regression coefficients of a hierarchical
spatio-temporal state-space model estimated by EM, and of its
spatially-clustered extension. The penalty leaves the E-step untouched, so the
algorithm remains an EM on the penalized observed-data likelihood; the M-step
for the coefficients keeps a closed form under a ridge and is solved exactly by
coordinate descent otherwise. Two features distinguish the problem from the
usual regularized regression. First, the relevant cross-product is a
generalized least squares one, carrying the inverse of an estimated spatial
covariance, so the scaling that makes the penalty comparable across covariates
is not the Euclidean standardization. Second, in the clustered model the
penalty operates within regimes whose size is itself a constraint of the
procedure, which is exactly where collinearity binds.
\end{abstract}

\section{Motivation}

Two things make a penalty on the regression coefficients worth having in this
class of models, and neither is the usual high-dimensional argument.

The first is that the covariates of an environmental monitoring network are
strongly collinear by construction. Meteorological drivers are measured on
overlapping supports, chemical species share sources, and the design matrix of a
regional network routinely has condition numbers in the hundreds. The variance
of the maximum likelihood estimator of $\beta$ is governed by the inverse of
$M = \sum_t X_t'\Se^{-1}X_t$, and it is that inverse that a ridge stabilizes.

The second is specific to the clustered model. A spatially-clustered STEM model
fits a complete state-space model inside every regime; the regimes are subsets
of the network, and their size is bounded below by the requirement that each of
them support its own variance components and range. Collinearity that is
tolerable on a network of two hundred stations is not tolerable on a regime of
twelve, and the binding constraint on how many regimes may be entertained is,
in practice, the smallest one that can still be fitted. A penalty that reduces
the effective dimension of the regression inside a regime relaxes that
constraint directly.

\section{The model}

For $t = 1,\dots,T$ and $d$ locations with coordinates $s_1,\dots,s_d$,
\begin{align}
  z_t &= X_t\beta + K y_t + e_t, & e_t &\sim N_d(0,\Se), \label{eq:meas}\\
  y_t &= G\,y_{t-1} + \eta_t,    & \eta_t &\sim N_p(0,\Sig_\eta),
  \label{eq:state}
\end{align}
with $X_t$ the $d\times r$ design of period $t$, $K$ a known $d\times p$
loading matrix, $y_0 \sim N_p(m_0, C_0)$, and
\begin{equation}
  \Se = \sigma^2_\varepsilon I_d + \sigma^2_\omega\, C(h;\theta),
  \qquad C(h;\theta) = \exp(-\theta h),
  \label{eq:cov}
\end{equation}
$h$ the matrix of distances. Write
$\psi = (\beta,\sigma^2_\varepsilon,\sigma^2_\omega,\theta,G,\Sig_\eta,m_0)$.

The spatially-clustered model partitions the locations into $k$ regimes
$I_1,\dots,I_k$ and gives each of them its own $\psi_g$, the partition being
estimated jointly with the parameters by maximizing a Potts-penalized
likelihood. Everything below applies within a regime, so we develop it for the
pooled model and return to the clustered case in Section~\ref{sec:scstem}.

\section{The penalized objective}

Let $\ell(\psi)$ be the observed-data log-likelihood of \eqref{eq:meas}. We
consider
\begin{equation}
  \ell_\pi(\psi) = \ell(\psi) - \pi(\beta),
  \qquad
  \pi(\beta) = \lambda\left\{
    \alpha\,\lVert D\beta\rVert_1
    + \tfrac{1-\alpha}{2}\,\beta'D\beta \right\},
  \label{eq:pen}
\end{equation}
with $\lambda \ge 0$, $\alpha\in[0,1]$ and $D = \mathrm{diag}(w_1,\dots,w_r)$,
$w_j\in\{0,1\}$, the indicator of the penalized coordinates. This is the
parameterization of \citet{ZouHastie2005}: $\alpha=0$ is a ridge, $\alpha=1$ a
lasso, and anything in between an elastic net.

Only $\beta$ is penalized. The variance components, the range, the transition
matrix and the initial state are left at their maximum likelihood values. That
is a modelling decision and not a technical one: $\sigma^2_\varepsilon$,
$\sigma^2_\omega$ and $\theta$ describe the error process rather than the mean,
and shrinking them would change the object being estimated rather than
stabilize its estimation.

\subsection{The intercept}

The default takes $w_1 = 0$ for the intercept, and here the usual convention
carries a substantive reason. The state equation \eqref{eq:state} has an
initial mean $m_0$ which is estimated, and the latent process it starts feeds
into \eqref{eq:meas} additively through $K$. When $K = \mathbf{1}_d$, as it is
in the clustered model within each regime, the intercept and the level of the
latent process are identified only through the dynamics: a constant may be moved
from $\beta_0$ into $m_0$ and back with a compensating change in the transient.
Shrinking $\beta_0$ therefore does not shrink ``the level''; it moves the level
into $m_0$, at a rate that depends on $G$. The penalty would be doing something
other than what it appears to do.

\section{EM with a penalty}

Let $Q(\psi\mid\psi^{(i)}) = E_{\psi^{(i)}}[\,\ell_c(\psi)\mid z\,]$ be the
usual expected complete-data log-likelihood, the expectation taken over the
latent states $y_{0:T}$.

\begin{proposition}[The E-step is unchanged and the algorithm is monotone]
\label{prop:monotone}
Define the penalized EM map by
\begin{equation}
  \psi^{(i+1)} = \argmax_{\psi}\;\bigl\{ Q(\psi\mid\psi^{(i)}) - \pi(\beta) \bigr\}.
  \label{eq:mstep}
\end{equation}
Then $\ell_\pi(\psi^{(i+1)}) \ge \ell_\pi(\psi^{(i)})$ for every $i$.
\end{proposition}

\begin{proof}
Write $\ell(\psi) = Q(\psi\mid\psi^{(i)}) - H(\psi\mid\psi^{(i)})$ with
$H(\psi\mid\psi^{(i)}) = E_{\psi^{(i)}}[\log f(y_{0:T}\mid z,\psi)\mid z]$, and
recall $H(\psi\mid\psi^{(i)}) \le H(\psi^{(i)}\mid\psi^{(i)})$ by Jensen's
inequality. Since $\pi$ is a function of $\beta$ alone and does not involve
$y_{0:T}$, it passes through the conditional expectation unchanged and
\[
  \ell_\pi(\psi) = \bigl\{Q(\psi\mid\psi^{(i)}) - \pi(\beta)\bigr\}
                   - H(\psi\mid\psi^{(i)}).
\]
The first brace is maximized at $\psi^{(i+1)}$ by construction and the second
term is minimized at $\psi^{(i)}$, so the difference cannot decrease.
\end{proof}

The proposition is the reason the whole construction is cheap: the Kalman
filter and smoother, which are the expensive part of the algorithm, are not
touched at all. Only the point returned by the M-step changes.

\begin{remark}
Proposition~\ref{prop:monotone} concerns the pooled model. In the clustered
model the objective also carries the partition, which is updated by iterated
conditional modes, and the procedure is not globally monotone for reasons
unrelated to the penalty.
\end{remark}

\section{The M-step for the coefficients}

Conditionally on the other blocks, the part of $Q$ that involves $\beta$ is
\begin{equation}
  Q_\beta(\beta) = -\tfrac{1}{2}\beta' M \beta + \beta' v + \mathrm{const},
  \qquad
  M = \sum_{t=1}^{T} X_t'\Se^{-1}X_t,
  \quad
  v = \sum_{t=1}^{T} X_t'\Se^{-1}\hat v_t,
  \label{eq:Mv}
\end{equation}
where $\hat v_t = \hat z_t - K\hat y_t$ is the completed observation of period
$t$ net of the smoothed latent contribution --- ``completed'' because under
missing data $\hat z_t$ carries the conditional expectation of what was not
observed, and the conditional variance is added back into the second moments
elsewhere. The M-step therefore maximizes
\begin{equation}
  -\tfrac{1}{2}\beta'M\beta + \beta'v
  - \lambda\alpha\lVert D\beta\rVert_1
  - \tfrac{\lambda(1-\alpha)}{2}\beta'D\beta.
  \label{eq:mstepbeta}
\end{equation}

\subsection{Ridge: a closed form}

At $\alpha = 0$ the objective is a smooth concave quadratic and
\begin{equation}
  \hat\beta = \bigl( M + \lambda D \bigr)^{-1} v.
  \label{eq:ridge}
\end{equation}
$M$ and $v$ are already assembled by the M-step, and $r$ is small, so the cost
of \eqref{eq:ridge} over the unpenalized $M^{-1}v$ is one addition on a
diagonal.

\subsection{Lasso and elastic net: coordinate descent}

For $\alpha > 0$ the objective \eqref{eq:mstepbeta} is concave and its
non-smooth part is separable across coordinates, which is the setting in which
cyclic coordinate descent converges to a global maximizer
\citep{Tseng2001,Friedman2010}. Fixing all coordinates but the $j$th, the
first-order condition is
\[
  v_j - \sum_{l \ne j} M_{jl}\beta_l - M_{jj}\beta_j
  - \lambda(1-\alpha)w_j\beta_j
  - \lambda\alpha w_j\,\partial|\beta_j| \ni 0,
\]
whose solution is the soft-thresholding
\begin{equation}
  \beta_j \;\leftarrow\;
  \frac{\mathcal{S}\!\left(v_j - \sum_{l\ne j}M_{jl}\beta_l,\;
        \lambda\alpha w_j\right)}
       {M_{jj} + \lambda(1-\alpha)w_j},
  \qquad
  \mathcal{S}(x,\tau) = \sgn(x)\,(|x|-\tau)_+ .
  \label{eq:cd}
\end{equation}
An unpenalized coordinate has $w_j = 0$ and \eqref{eq:cd} reduces to its
ordinary conditional least squares update. Because the maximizer of the M-step
is attained exactly, \eqref{eq:mstep} is a genuine EM step and not merely a
generalized one; the cost is $O(r^2)$ per sweep on a matrix that is already
formed.

\subsection{The metric, and why the usual standardization is wrong here}
\label{sec:metric}

$M$ in \eqref{eq:Mv} is a \emph{generalized} least squares cross-product: it
carries $\Se^{-1}$, an estimated spatial covariance that changes from one EM
iteration to the next. The effective design is therefore
$\Se^{-1/2}X_t$ and not $X_t$, and standardizing the columns of $X$ to unit
Euclidean norm --- the default of every lasso implementation --- does not make
$\lambda$ comparable across covariates.

The scaling that does is the one that puts the diagonal of $M$ at one. Writing
$s_j = M_{jj}^{1/2}$, $S = \mathrm{diag}(s)$, $\tilde M = S^{-1}MS^{-1}$ and
$\tilde v = S^{-1}v$, the penalty is applied to $\tilde\beta = S\beta$ and the
solution returned on the original scale. Under this convention $\lambda$ is
measured in units of the GLS information each covariate carries, which is what
makes a single $\lambda$ meaningful across covariates and, in the clustered
model, across regimes of different sizes and different error covariances.

\begin{remark}
$s_j$ depends on $\Se$ and therefore changes between EM iterations. The scaling
is recomputed at each M-step. Since it is undone before the coefficients are
returned, the sequence of iterates is well defined; what is not constant across
iterations is the geometry in which a fixed $\lambda$ acts, and this is the
price of applying a penalty inside an algorithm that also estimates the metric.
\end{remark}

\section{Effective degrees of freedom and the information criteria}

The information criteria used to select $k$ and the Potts penalty count free
parameters. With a penalty in force the nominal count $r$ overstates the
flexibility of the fit, and the effective number must replace it.

For the ridge, the fitted values are linear in the response with hat matrix
$H_\lambda = X(M+\lambda D)^{-1}X'\Se^{-1}$ in the GLS metric, so
\begin{equation}
  \mathrm{df}(\lambda) = \tr\!\left\{ M\,(M+\lambda D)^{-1} \right\},
  \label{eq:dfridge}
\end{equation}
which decreases from $r$ at $\lambda = 0$ to the number of unpenalized
coordinates as $\lambda\to\infty$.

For the lasso, the number of non-zero coefficients is unbiased for the degrees
of freedom \citep{ZouHastieTibshirani2007}; for the elastic net the
corresponding quantity is \eqref{eq:dfridge} restricted to the active set
$\mathcal{A} = \{j : \hat\beta_j \ne 0\}$, with the ridge part of the penalty
only:
\begin{equation}
  \mathrm{df}(\lambda,\alpha) =
  \tr\!\left\{ M_{\mathcal{A}}\,
    \bigl(M_{\mathcal{A}} + \lambda(1-\alpha)D_{\mathcal{A}}\bigr)^{-1}\right\}.
  \label{eq:dfen}
\end{equation}
The criteria then read, for a fit with $k$ regimes,
\[
  \mathrm{BIC} = -2\hat\ell + \log(N)\Bigl\{
    \textstyle\sum_{g=1}^{k}\mathrm{df}_g(\lambda,\alpha) + k\,q \Bigr\},
\]
with $q$ the number of parameters of a regime other than the coefficients and
$N$ the number of observations, and analogously for AIC and KIC.

\begin{remark}
$\hat\ell$ is the \emph{unpenalized} log-likelihood evaluated at the penalized
estimate. Using the penalized objective in its place would double-count the
penalty, once through the objective and once through the degrees of freedom.
\end{remark}

\section{Choosing $\lambda$ and $\alpha$}

$(\lambda,\alpha)$ join $(k,\varphi)$ as hyperparameters and are selected the
same way. Two routes are available and they do not agree in general.

\paragraph{Information criteria.} Cheap, since the grid is already being
traversed for $(k,\varphi)$, and consistent with the rest of the selection
apparatus once \eqref{eq:dfridge}--\eqref{eq:dfen} are used. The known
difficulty is that in-sample criteria on a model whose partition is itself
estimated are optimistic, and the optimism grows with $k$.

\paragraph{Spatio-temporal cross-validation.} The honest route, and the one we
recommend when the computational budget allows it. The blocking scheme matters:
random $K$-fold over the $(t,i)$ cells is anti-conservative here because the
latent process and the spatial field make neighbouring cells nearly duplicates.
Leave-$L$-locations-out, leave-$H$-times-out and the combined
leave-$L$-locations-and-$H$-times-out schemes address the two dependencies
separately and jointly; see \citet{OttoFassoMaranzano2024} for the review and
\citet{Meyer2018} for the implementation these follow. The predictive quantity
to score is the fitted value at held-out cells, which the package computes
through the conditional expectation given the observed sub-vector.

\section{Inference}

The parametric bootstrap already available for the model generates a replicate
from the fitted parameters and re-estimates. With a penalty two changes are
required.

First, if $(\lambda,\alpha)$ were selected from the data, the selection must be
repeated inside every bootstrap draw. Otherwise the intervals condition on a
tuning parameter that is itself an estimate, in exactly the way that intervals
conditioning on an estimated partition understate the uncertainty of the
clustered model.

Second, when $\alpha > 0$ the estimator is not smooth and the sampling
distribution of a coefficient at the boundary is not asymptotically normal; the
percentile and bias-corrected intervals remain usable as descriptions of the
bootstrap distribution but do not have their usual coverage interpretation.
Post-selection inference for the active set requires the conditional machinery
of \citet{LeeSunSunTaylor2016}, which we do not develop here. For a ridge,
$\alpha = 0$, the estimator is smooth and this difficulty does not arise,
though the usual bias remains.

\section{The clustered model}
\label{sec:scstem}

Within a regime, everything above applies verbatim with $M_g$, $v_g$ and
$\Se^{(g)}$ in place of their pooled counterparts. Three points are specific to
the clustered case.

\paragraph{A common $\lambda$ across regimes is meaningful only after scaling.}
Regimes differ in size and in error covariance, so $M_g$ differs in magnitude
between them. The scaling of Section~\ref{sec:metric} is what makes a single
$(\lambda,\alpha)$ interpretable across the partition; without it the penalty
would bind hardest on the smallest regime for a reason that has nothing to do
with its collinearity.

\paragraph{The penalty interacts with the minimum regime size.} The lower bound
on $n_g$ exists because a regime must support its own parameters. A penalty
reduces the effective number of coefficients, and with it the bound. This is the
most substantive consequence for practice: it is what would allow finer
partitions to be entertained on a network of a given size.

\paragraph{A fused alternative.} The penalty considered here shrinks each
$\beta_g$ towards zero independently. An alternative, arguably better matched to
the clustered model, shrinks them towards each other,
\[
  \pi_{\mathrm{fused}}(\beta_1,\dots,\beta_k)
  = \lambda \sum_{g < g'} \nu_{gg'} \lVert D(\beta_g - \beta_{g'})\rVert,
\]
with $\nu_{gg'}$ a weight, possibly spatial. That penalty addresses the
smallness of the regimes rather than the collinearity of the covariates, it
couples the M-steps of different regimes, and it makes the number of distinct
coefficient vectors an outcome of the fit rather than a choice. It is a
different paper.

\section{Algorithm}

\begin{quote}\small
\textbf{input} data $z$, design $X$, coordinates, starting $\psi^{(0)}$,
$(\alpha,\lambda)$, weights $w$\\[2pt]
\textbf{repeat} for $i = 0,1,2,\dots$\\
\hspace*{1em} 1. E-step: Kalman filter and smoother at $\psi^{(i)}$; form
$\hat y_t$, $C^s_t$, $\hat z_t$\\
\hspace*{1em} 2. M-step, other blocks: closed forms for $\sigma^2_\omega$,
$G$, $\Sig_\eta$, $m_0$; Newton--Raphson for $(\theta,\sigma^2_\varepsilon)$\\
\hspace*{1em} 3. M-step, coefficients:\\
\hspace*{2em} (a) assemble $M$, $v$ as in \eqref{eq:Mv}\\
\hspace*{2em} (b) $s_j \leftarrow M_{jj}^{1/2}$;
$\tilde M \leftarrow S^{-1}MS^{-1}$, $\tilde v \leftarrow S^{-1}v$\\
\hspace*{2em} (c) if $\alpha = 0$: $\tilde\beta \leftarrow
(\tilde M + \lambda D)^{-1}\tilde v$\\
\hspace*{2em} \phantom{(c)} else: cycle \eqref{eq:cd} over $j$ until
convergence\\
\hspace*{2em} (d) $\beta^{(i+1)} \leftarrow S^{-1}\tilde\beta$;
record $\mathrm{df}$ from \eqref{eq:dfridge} or \eqref{eq:dfen}\\
\textbf{until} the relative changes in $\psi$ and in $\ell_\pi$ are both below
the tolerance
\end{quote}

The stopping rule needs one adjustment when $\alpha > 0$: the relative criterion
on the parameter vector is undefined at a coordinate that has just become
exactly zero, and the active set may change from one iteration to the next
without the objective moving appreciably. The implementation floors the
denominator and tests the objective as well as the parameters.

\section{What remains to be done}

\begin{enumerate}
\item \textbf{Simulation evidence.} The design of the companion software paper
provides the machinery: what is missing is a factor for the collinearity of the
covariates, which is the quantity the penalty is supposed to address and which
that design holds fixed at a single covariate.
\item \textbf{Consistency and rates.} Under what conditions on $T$, $d$, and the
decay of the spatial correlation is $\hat\beta(\lambda)$ consistent, and at what
rate? The effective sample size is not $Td$: the spatial field and the latent
process both reduce it, and the reduction depends on $\theta$ and $G$.
\item \textbf{The degrees of freedom of \eqref{eq:dfen}} are the fixed-design
ones. Here the metric is estimated, and whether the correction matters at
realistic $T$ is an open question we have not examined.
\item \textbf{Selection consistency} of the lasso part in the presence of an
estimated spatial covariance, and whether the irrepresentable condition can be
stated in the GLS metric.
\item \textbf{The fused penalty} of Section~\ref{sec:scstem}, which is a
separate development and probably the more interesting one.
\end{enumerate}

\begin{thebibliography}{9}
\small
\bibitem[Friedman et~al.(2010)]{Friedman2010}
Friedman, J., Hastie, T., Tibshirani, R. (2010)
Regularization paths for generalized linear models via coordinate descent.
\emph{Journal of Statistical Software}, 33, 1--22.

\bibitem[Lee et~al.(2016)]{LeeSunSunTaylor2016}
Lee, J.~D., Sun, D.~L., Sun, Y., Taylor, J.~E. (2016)
Exact post-selection inference, with application to the lasso.
\emph{The Annals of Statistics}, 44, 907--927.

\bibitem[Meyer et~al.(2018)]{Meyer2018}
Meyer, H., Reudenbach, C., Hengl, T., Katurji, M., Nauss, T. (2018)
Improving performance of spatio-temporal machine learning models using forward
feature selection and target-oriented validation.
\emph{Environmental Modelling and Software}, 101, 1--9.

\bibitem[Otto et~al.(2024)]{OttoFassoMaranzano2024}
Otto, P., Fasso, A., Maranzano, P. (2024)
A review of regularised estimation methods and cross-validation in
spatiotemporal statistics.
\emph{Statistics Surveys}, 18, 299--340.

\bibitem[Tseng(2001)]{Tseng2001}
Tseng, P. (2001)
Convergence of a block coordinate descent method for nondifferentiable
minimization.
\emph{Journal of Optimization Theory and Applications}, 109, 475--494.

\bibitem[Zou and Hastie(2005)]{ZouHastie2005}
Zou, H., Hastie, T. (2005)
Regularization and variable selection via the elastic net.
\emph{Journal of the Royal Statistical Society B}, 67, 301--320.

\bibitem[Zou et~al.(2007)]{ZouHastieTibshirani2007}
Zou, H., Hastie, T., Tibshirani, R. (2007)
On the degrees of freedom of the lasso.
\emph{The Annals of Statistics}, 35, 2173--2192.
\end{thebibliography}

\end{document}
